problem:  
Let \(V^n \subset S^{n+k}\) be a stationary integral \(n\)-varifold whose regular part is a smooth minimal submanifold.  
For each singular point \(p \in \mathrm{sing}\,V\), assume there exists a **unique tangent cone**
\[
C_p = \mathbb{R}_+ \times L_p,
\]
where \(L_p \subset S^{n+k-1}\) is a smooth minimal \((n-1)\)-submanifold — the **link** of the cone.

Assume that each link \(L_p\) is **unstable**, but with uniformly bounded **Morse index**:
\[
\mathrm{ind}(L_p) \le N,
\]
for some fixed integer \(N\).  (So each \(\mathcal{J}_{L_p}\), the Jacobi operator on the normal bundle, has at most \(N\) negative eigenvalues.)

**Question:**  
Under this assumption of *uniform index bound* on all tangent cone links, must the family
\[
\Sigma(V) = \{[L_p] : p \in \mathrm{sing}\,V\}
\]
be **finite up to smooth deformation and ambient isometries**?  

Equivalently, does there exist, for given \(n,k,N\), a finite set of smooth minimal submanifolds  
\(\{L^{(1)}, \ldots, L^{(m)}\} \subset S^{n+k-1}\)  
such that every tangent cone link \(L_p\) of \(V\) is smoothly deformable (via minimal embeddings in the same sphere) to one of the \(L^{(i)}\)?  

If not, can one construct a stationary minimal varifold in \(S^{n+k}\) possessing an **infinite discrete family of non-isometric tangent cone links**, all having uniformly bounded Morse index?  

That is, can index bounds alone prevent or fail to prevent the proliferation of distinct local models of singularities in positively curved ambient spaces?

Why is it a "good" problem:  
This problem corrects the logical inconsistency present in earlier formulations: it removes the incompatible pairing of stability (which forbids negative eigenvalues) with a bound on negative spectrum. The resulting question lies at a precise and fertile intersection of **spectral geometry**, **moduli theory of minimal submanifolds in spheres**, and **regularity theory for minimal varifolds**.  

The hypothesis — uniformly bounded index for tangent cone links — is a natural analytic control condition and one studied intensely in compactness and curvature estimates for minimal submanifolds. Yet, at present, no theorem asserts or rules out finiteness (up to smooth deformation) of minimal submanifolds of spheres under such uniform index bounds.  

A positive answer would unveil a **spectral–geometric finiteness principle**: bounded instability implies only finitely many possible link geometries for singularities of stationary varifolds in \(S^{n+k}\). That would tie analytic spectral control directly to global geometric classification — a major conceptual advance connecting *index theory* with *singular moduli rigidity*.  
A negative answer, via construction of infinitely many distinct link types satisfying the same index bound, would show that spectral data alone are insufficient to regulate geometric moduli in positive curvature, exposing the limits of spectral control in geometric measure theory.  

The statement is simple — bounded index implies finiteness of geometric link types — yet its resolution would require breakthroughs merging GMT regularity, minimal surface compactness, and analytic index theory. It thus embodies the “narrow entrance, profound depth” quality of a genuinely good research problem.